\documentclass[aspectratio=169,10pt]{beamer}
\usepackage[utf8]{inputenc}
\usepackage{amsmath}
\usetheme{default}
\usecolortheme{dove}
\setbeamertemplate{navigation symbols}{}
\setbeamercolor{structure}{fg=black}
\setbeamertemplate{itemize item}{--}
\setbeamertemplate{frametitle}{\vspace{2mm}\textbf{\insertframetitle}\vspace{1mm}}
\setbeamersize{text margin left=9mm, text margin right=9mm}

\begin{document}
\begin{frame}
\frametitle{Empirical strategy --- Gillmore (2025)}
{\footnotesize Difference-in-differences: municipal seismic intensity $\times$
birth-cohort exposure to the 2010 Chilean earthquake (27-F).}

\vspace{2mm}
\begin{equation*}
Y_{ijtw}=\beta\,\bigl(\mathit{Affected}_{it}\times \mathit{Mercalli}_{j}\bigr)
+\gamma\,\mathit{Affected}_{it}+X_{ijt}+\psi_{j}+\theta_{t}+\lambda_{jt}
+\eta_{w}+e_{ijtw}
\tag{2.1}
\end{equation*}

\vspace{1mm}
\begin{itemize}\setlength\itemsep{1.5pt}\footnotesize
  \item $i$ child, $j$ municipality, $t$ birth cohort, $w$ survey wave (2012, 2017)
  \item $\mathit{Affected}_{it}=1$ if conceived before 27-F: in utero to age~5
        at impact (vs.\ cohorts conceived after)
  \item $\mathit{Mercalli}_{j}$: modified Mercalli intensity of the municipality
        (Astroza et al.\ 2010); robustness: USGS peak ground acceleration
  \item $X_{ijt}$: child, mother and household controls (incl.\ pre-quake);
        $\psi_{j},\theta_{t},\eta_{w}$: municipality, cohort and wave fixed
        effects; $\lambda_{jt}$: municipality linear trends
\end{itemize}

\vspace{2mm}
{\footnotesize \textbf{Identification.} $\beta$ = intention-to-treat effect of one
additional Mercalli unit on exposed cohorts. Parallel trends conditional on
$\lambda_{jt}$; post-quake conceptions possibly indirectly exposed
$\Rightarrow$ $\beta$ is a lower bound. Standard errors clustered by
municipality.}

\vspace{1.5mm}
{\footnotesize \textbf{Validity.} Placebo on post-quake cohorts: null. Migration
between waves: 0.78\% (affected$\rightarrow$unaffected 0.02\%). Fertility,
planned births and sex ratio: null in vital records. Attrition uncorrelated
with intensity.}

\vfill
{\tiny Eq.\ (2.1): Gillmore, \emph{Natural Disasters and Early Child
Development: Evidence from an Earthquake} --- UT Austin dissertation (2023,
ch.~2); SSRN WP 5106675 (2025); \emph{Economics of Education Review} (2026).
Published version: exposure prenatal--age 4. Outcomes: TVIP/Peabody, Battelle,
CBCL.}
\end{frame}
\end{document}
